% SIGMOD 2027 研究论文投稿。格式依据官方 CFP：
%   https://2027.sigmod.org/calls_papers_sigmod_research.shtml
%   - 正文 ≤ 12 页，参考文献不计入且页数不限（录用后 camera-ready 为 13 页且不得有附录）
%   - 双盲：正文任何位置不得出现作者、单位、致谢、资助来源
%   - 2 栏 ACM 格式，字号/页边距/栏间距/行距一律不得改动
%   - 提交 PDF ≤ 10 MB，Letter 纸
%
% 类选项：sigconf = ACM 会议 2 栏格式；review = 给正文加行号方便审稿人引用；
% anonymous = 隐去作者与单位。定稿时把 review,anonymous 两项删掉即可。
\documentclass[sigconf,review,anonymous]{acmart}

% —— 投稿期去掉版权块 ——
% acmart 默认会在首页脚印 ACM 版权声明和 DOI，那些数字要等录用后从 ACM eRights
% 系统拿。投稿阶段用下面三行清掉；这只影响首页脚注区，不改动任何版面度量，
% 不违反 CFP 的"格式不得改动"。录用后连同 \acmConference 一起换成 eRights 给的内容。
\setcopyright{none}
\settopmatter{printacmref=false}
\renewcommand\footnotetextcopyrightpermission[1]{}

% —— camera-ready 时启用，内容以 ACM eRights 邮件为准 ——
% \copyrightyear{2027}
% \acmYear{2027}
% \setcopyright{acmlicensed}
% \acmConference[SIGMOD '27]{International Conference on Management of Data}{May 31--June 4, 2027}{Huntington Beach, CA, USA}
% \acmDOI{10.1145/nnnnnnn.nnnnnnn}
% \acmISBN{978-x-xxxx-xxxx-x/YY/MM}

\begin{document}

\title{Collaborating without Coordinating: Data Management for Knowledge Accumulated by LLM Agents}

% anonymous 选项会自动把下面这些替换成 "Anonymous Author(s)"，所以真名可以先写着，
% 不必等到定稿再填。但正文、致谢、图表、URL、自引措辞里的身份线索必须自己管。
\author{Author One}
\affiliation{%
  \institution{Institution}
  \city{City}
  \country{Country}}
\email{one@example.com}

\renewcommand{\shortauthors}{Anonymous et al.}

\begin{abstract}
LLM agents now both use and produce knowledge: as they read, compare, and
verify, they can build on what earlier agents recorded and leave records for
later ones. When many agents share one long-lived store, no single agent's word
is authoritative: agents may work without coordinating or legitimately
disagree, and the sources they draw on change. We argue that such a store
cannot decide who is right, but must keep every write accountable: converging
what rules can decide, preserving disagreement with its grounds, and exposing
what may be stale. We study this problem in scientific literature, where it is
acute and checkable, and present PaperWeave, a data management system whose
data model separates source-grounded records, revised with full history, from
interpretations that coexist with their provenance. It adapts classical entity
resolution, provenance, view invalidation, and data versioning to agent writers,
with contracts for self-describing contributions and retrieval interfaces for
readers who were not present when those contributions were written. We frame
the evaluation around an intent-grounded workload to test whether these
mechanisms limit degradation as writers and updates accumulate, and whether
construction costs amortize across later readers.
% TODO: Replace the evaluation scope above with measured findings once experiments are complete.
\end{abstract}

% CCS 概念与关键词：投稿时非必需，camera-ready 必填。
% CCS 代码到 https://dl.acm.org/ccs 用 CCS 工具生成后粘贴回来。
% \begin{CCSXML} ... \end{CCSXML}
% \ccsdesc[500]{Information systems~Information extraction}
% \keywords{keyword one, keyword two}

\maketitle

\section{Introduction}
\label{sec:intro}


\section{Conclusion}
\label{sec:conclusion}

% 页数闸门的锚点：Makefile 的 `make check` 从 .aux 里读这个 label 所在页码，
% 判断正文是否超过 12 页。它必须紧挨在参考文献之前。
\label{endofbody}

% 双盲期不得出现致谢。录用后再取消注释。
% \begin{acks}
% Acknowledgements go here.
% \end{acks}

\bibliographystyle{ACM-Reference-Format}
\bibliography{refs}

\end{document}
